\section{Certificate-Checked Passes}
\label{sec:checked}
The passes of Section~\ref{sec:post} are local rewrites, and their proofs follow the rewrite. Two further reductions need a global analysis that is hard to prove correct and easy to check: the removal of transitions that another transition dominates, and the reuse of clocks. For both, an unverified OCaml procedure proposes a \emph{certificate}, a boolean checker extracted from Rocq tests local conditions on it, and the pass returns the automaton unchanged when the check fails. The Rocq theorem is that the language is preserved when the check succeeds, so it quantifies over every certificate, including wrong ones, and nothing is assumed about the procedure. The scheme is the one of the validation of register allocation in CompCert~\cite{rideau_leroy_cc2010}. Both passes run after the optimization rounds and before export; the pruning is preceded by a verified splitting of transitions.

\subsection{Removal of dominated transitions}
\label{sec:prune}
Two transitions out of one location may overlap: whenever the second is enabled so is the first, and the first leads to a location that can do everything that the target of the second can do. The second is then redundant for the language but not for a model checker, which explores both. Syntactic subsumption (Section~\ref{sec:post}) sees only pairs with the same target and resets; simulation-based reductions of B\"uchi automata~\cite{clemente_mayr_lmcs19} are untimed, and clock guards and resets must be related in a timed setting.

For a set $K\subseteq E$ of transitions to keep, the pruned automaton is $\mathcal A|_K=(L,\ell_0,X,K,\mathit{Acc})$. Let $\sim$ be a relation on $L$ ($q_1\sim q_2$: $q_1$ simulates $q_2$). A transition $t_1=(q_1,\alpha_1,g_1,Z_1,q_1')$ \emph{dominates} $t_2=(q_2,\alpha_2,g_2,Z_2,q_2')$ on a cell $c$ of letters (see below) if both labels hold on $c$, $g_2$ implies $g_1$, $Z_1$ and $Z_2$ agree on the clocks still read after the step, and $q_1'\sim q_2'$. The certificate is $(K,\sim)$ with, for each location, the clocks whose values in the two runs must be equal ($\mathit{live}$) or ordered ($\mathit{lle}$, enough for upper-bound guards). The checker \path{prune_ok} tests that $\ell_0\sim\ell_0$, that $q_1\sim q_2$ and $q_2$ accepting imply $q_1$ accepting, and that for $q_1\sim q_2$ every transition out of $q_2$ is dominated, on every cell where it is enabled, by a kept transition out of $q_1$; the exact conditions on clocks are in Supplement Section~\ref{S-sec:prune-detail}.

\begin{theorem}[Pruning]
\label{thm:prune}
For every certificate producer and every TBA $\mathcal A$, the automaton returned by \path{prune_pass} accepts exactly the timed words accepted by $\mathcal A$.
\end{theorem}
\begin{proof}[Idea]
Every run of the pruned automaton is a run of $\mathcal A$. Conversely, a run of $\mathcal A$ is mapped position by position to a run of the pruned automaton: at each step the transition taken is replaced by the kept transition that dominates it, and an invariant states that the two locations are related by $\sim$ and that the valuations agree on the live clocks and are ordered on the clocks of $\mathit{lle}$. Acceptance is preserved because accepting positions map to accepting positions. The mechanized result is \path{prune_pass_accepts}; if the check fails, the pass is the identity.
\end{proof}

\paragraph{Splitting by cells of letters.} A transition whose label holds on several letters may be dominated on some of them only. \path{split_cells} replaces each transition by one copy for every \emph{cell} of letters in its label, one cell for each action mentioned in a label of the automaton and one for all other actions, each copy labelled by the literals that characterize its cell. This makes the label test exact, and \path{split_cells_accepts} shows that the language is unchanged. Splitting alone enlarges the automaton; it pays when the pruning that follows removes the copies that other transitions dominate (Section~\ref{sec:eval-checked}).

\subsection{Reuse of clocks by liveness}
\label{sec:reuse}
Translators introduce one clock for each timed subformula, and clocks that are never read at the same time can share one variable, which lowers the dimension of every zone. Liveness, interference graphs and colouring are the method of Saeedloei and Kluźniak for clock allocation in timed automata~\cite{saeedloei_kluzniak_hscc2018}; ours is a Rocq-checked version for the semantics of this chain. The certificate is an allocation $\alpha:\mathcal C\to\mathcal C$ and sets $\mathit{live}(q)$ that over-approximate the clocks read before being reset from $q$. The renamed automaton is $\alpha(\mathcal A)=(L,\ell_0,\alpha(X),\{(\ell,a,\alpha(g),\alpha(Z),\ell')\mid(\ell,a,g,Z,\ell')\in E\},\mathit{Acc})$. The checker \path{alloc_ok} tests, for every transition $t$ from $q$ to $q'$: \textbf{(G)} the clocks of its guard are in $\mathit{live}(q)$; \textbf{(F)} $\mathit{live}(q')\subseteq\mathit{reset}(t)\cup\mathit{live}(q)$; \textbf{(R)} if $r$ is reset and $y\in\mathit{live}(q')$ with $\alpha(r)=\alpha(y)$ then $r=y$; and \textbf{(I)} $\alpha$ is injective on $\mathit{live}$ of the initial location.

\begin{theorem}[Clock reuse]
\label{thm:reuse}
If $(\alpha,\mathit{live})$ passes \path{alloc_ok} for $\mathcal A$, then $\alpha(\mathcal A)$ accepts exactly the timed words accepted by $\mathcal A$; \path{alloc_pass} applies $\alpha$ only after a successful check.
\end{theorem}
\begin{proof}[Idea]
Along a run, a class of $\alpha$ is reset when one of its members is, and every live clock has the value of its class (by (I) at the start, by (F) and (R) at each step); with (G) the guards of the renamed transition hold exactly when those of the original do, and the converse direction picks the first transition that maps to the one taken, so no choice axiom is used. The mechanized result is \path{try_allocate_accepts}; the conditions are discussed in Supplement Section~\ref{S-sec:alloc-detail}.
\end{proof}

The certifier computes the least liveness sets by a fixpoint, builds the interference graph (an edge between a reset clock and each clock live at the target, and between the clocks live initially) and colours it with DSATUR~\cite{brelaz_cacm79}; it is heuristic and need not find the fewest clocks.

\paragraph{Composition.} The checked passes compose with the translation: for every pruning certifier, every allocation certifier, every $n$ and every well-formed $f$, a timed word satisfies $f$ iff the export of $\mathsf{alloc\_pass}(\mathsf{prune\_pass}(\mathsf{split\_cells}(\optimize(n,\mathsf{compile}(f)))))$ accepts it (\path{MTL_to_exported_checked_correct}); the variant for the Spot-parametrized chain (\path{MTL_to_exported_checked_correct_with}) carries the same hypothesis on the propositional automaton as Theorem~\ref{thm:composition}. Under UPPAAL's initialization the guarantee needs the check of Section~\ref{sec:toolchain} on the final automaton: if \path{init_free} succeeds on the export of the checked chain, then $\mathsf{msat}(w,0,f)$ holds iff the exported automaton accepts $w$ with every clock initialized to 0 (\path{MTL_to_exported_checked_correct0_recur_weak_with}, and \path{MTL_to_exported_checked_correct0_with}), under the same Spot contract and for every certificate; the tool runs the check after pruning and clock reuse. \texttt{Print Assumptions} lists only the library axioms of Section~\ref{sec:toolchain}.
